\documentclass[12pt,a4paper]{article}
\usepackage[utf8]{inputenc}
\usepackage[T5]{fontenc}
\usepackage[vietnamese]{babel}
\usepackage[margin=1.5cm,top=1.6cm,bottom=1.5cm]{geometry}
\usepackage{amsmath,amssymb,array,tabularx,enumitem,fancyhdr,tikz,xcolor}
\usetikzlibrary{calc,angles,quotes,arrows.meta,babel}
\definecolor{ink}{RGB}{30,55,70}
\newcommand{\headerQuote}{%
  \ifcase\value{page}\relax
    \or Math gives us hope that every problem has a solution% Trang 1
    \or Math + You = Champion% Trang 2
    \or Where Frustration Becomes Inspiration% Trang 3
    \or Turning Confusion into Clarity% Trang 4
    \or Don't be an X, be a Y of mathematics% Trang 5
    \or Don't get mad. Learn to add% Trang 6
    \or Math gives us hope that every problem has a solution% Trang 7
    \or Math + You = Champion% Trang 8
    \or Where Frustration Becomes Inspiration% Trang 9
    \or Turning Confusion into Clarity% Trang 10
    \or Don't be an X, be a Y of mathematics% Trang 11
    \or Don't get mad. Learn to add% Trang 12
    \else Math gives us hope that every problem has a solution%
  \fi
}
\pagestyle{fancy}\fancyhf{}
\lhead{\small\textit{\headerQuote}}
\rhead{\small\textit{Người biên soạn: Dương Gia Bảo}}
\rfoot{\small Trang \thepage}
\renewcommand{\headrulewidth}{0.4pt}
\setlength{\headheight}{18pt}
\setlength{\parindent}{0pt}
\setlength{\parskip}{2pt}
\setlength{\tabcolsep}{5pt}
\renewcommand{\arraystretch}{1.16}
\newcommand{\parttitle}[1]{\par\vspace{3pt}{\large\bfseries\color{ink}#1}\par\vspace{2pt}}
\newcommand{\four}[4]{\par\noindent\begin{tabularx}{\linewidth}{@{}*{4}{>{\raggedright\arraybackslash}X}@{}}\textbf{A.} #1&\textbf{B.} #2&\textbf{C.} #3&\textbf{D.} #4\end{tabularx}}
\newcommand{\two}[4]{\par\noindent\begin{tabularx}{\linewidth}{@{}*{2}{>{\raggedright\arraybackslash}X}@{}}\textbf{A.} #1&\textbf{B.} #2\\[3pt]\textbf{C.} #3&\textbf{D.} #4\end{tabularx}}
\newcommand{\answerbox}{\hfill\textit{Đáp số:}\ \fbox{\rule{0pt}{13pt}\hspace{2.2cm}}}
\setlist[enumerate,1]{label=\textbf{Câu \arabic*.},leftmargin=1.5cm,labelsep=.18cm,itemsep=3.5pt,topsep=2pt,parsep=0pt}
\setlist[enumerate,2]{label=\alph*),leftmargin=.65cm,itemsep=2.5pt,topsep=2pt,parsep=0pt}
\begin{document}
\noindent
\begin{minipage}[c]{0.38\textwidth}
    \centering
    \textbf{LÀM TOÁN ĐỂ HẠNH PHÚC}\\[-2pt]
    \rule{3.2cm}{0.6pt}\\[4pt]
    \textbf{ĐỀ CHÍNH THỨC}
\end{minipage}\hfill
\begin{minipage}[c]{0.35\textwidth}
    \centering
    \textbf{ĐỀ ÔN THI GK1}\\
    \textbf{NĂM HỌC 2026 -- 2027}\\[2pt]
    \textbf{ĐỀ SỐ 01}\\
    \textbf{Môn: TOÁN -- Khối lớp 10}\\
    \textit{\small Thời gian làm bài: 90 phút\\(không kể thời gian phát đề)}
\end{minipage}

\vspace{0.2cm}
\noindent
\begin{minipage}[b]{0.78\textwidth}
    \textbf{Họ và tên:}\dotfill\textbf{SBD:}\makebox[3.5cm]{\dotfill}
\end{minipage}\hfill
\begin{minipage}[b]{0.20\textwidth}
    \hfill
    \setlength{\fboxrule}{0.8pt}%
    \fbox{\begin{minipage}{2.3cm}
        \centering
        \textbf{\small Mã đề thi}\\
        \textbf{0001}
    \end{minipage}}
\end{minipage}
\vspace{0.2cm}

\parttitle{PHẦN I. Trắc nghiệm nhiều lựa chọn}
\textit{Trả lời câu 1 đến câu 12. Mỗi câu chỉ chọn một phương án.}
\begin{enumerate}
\item Cho $90^\circ<\alpha<180^\circ$. Khẳng định nào sau đây đúng?
\four{$\sin\alpha<0$.}{$\cos\alpha<0$.}{$\tan\alpha>0$.}{$\cot\alpha>0$.}
\item Câu nào sau đây là một mệnh đề toán học?
\two{Bạn đã làm bài tập chưa?}{Hãy tính tổng hai số này!}{Bài toán này thật thú vị!}{$17$ là một số nguyên tố.}
\item Bất phương trình nào sau đây là bất phương trình bậc nhất hai ẩn $x,y$?
\four{$2x-3y\leq5$.}{$x^2+y>1$.}{$xy\leq4$.}{$x+y^2\geq0$.}
\item Cho $A=\{1;2;3\}$. Khẳng định nào sau đây đúng?
\four{$1\subset A$.}{$\{1\}\in A$.}{$\{1\}\subset A$.}{$4\in A$.}
\item Trong tam giác $ABC$, đặt $a=BC$, $b=CA$, $c=AB$. Công thức nào đúng?
\two{$a^2=b^2+c^2-2bc\cos A$.}{$a^2=b^2+c^2+2bc\cos A$.}{$a^2=b^2+c^2-2bc\sin A$.}{$a^2=b^2+c^2-2ab\cos A$.}
\item Cặp số $(x;y)$ nào sau đây là một nghiệm của $x+2y>4$?
\four{$(0;2)$.}{$(2;1)$.}{$(1;1)$.}{$(1;2)$.}
\item Câu nào sau đây là một mệnh đề chứa biến $x\in\mathbb R$?
\four{$2+3=5$.}{$x+2>5$.}{$9$ là số lẻ.}{$4<1$.}
\item Giá trị của $\sin150^\circ$ bằng
\four{$-\dfrac12$.}{$\dfrac{\sqrt3}{2}$.}{$\dfrac12$.}{$-\dfrac{\sqrt3}{2}$.}
\item Hệ nào sau đây là hệ bất phương trình bậc nhất hai ẩn $x,y$?
\two{$\begin{cases}x^2+y\leq4,\\x-y>1\end{cases}$}{$\begin{cases}x+y\leq4,\\2x-y>1\end{cases}$}{$\begin{cases}xy\leq4,\\x+y>1\end{cases}$}{$\begin{cases}x+y\leq4,\\x+y^2>1\end{cases}$}
\item Số tập con của tập hợp $M=\{2;5;8\}$ là
\four{$8$.}{$3$.}{$6$.}{$9$.}
\item Cho tam giác $ABC$ có $BC=a$, $CA=b$, $AB=c$ và diện tích $S$. Công thức nào đúng?
\four{$S=\dfrac12ab\cos C$.}{$S=ab\sin C$.}{$S=\dfrac12bc\sin C$.}{$S=\dfrac12ab\sin C$.}
\item Điểm nào thuộc miền nghiệm của hệ $\begin{cases}x\geq0,\ y\geq0,\\x+y\leq4,\\x+2y\geq4\end{cases}$?
\four{$(-1;3)$.}{$(3;2)$.}{$(1;2)$.}{$(0;1)$.}
\end{enumerate}
\newpage
\parttitle{PHẦN II. Trắc nghiệm đúng/sai}
\textit{Với mỗi ý a), b), c), d) trong từng câu, hãy xác định đúng hoặc sai.}
\begin{enumerate}
\item Một gian hàng chuẩn bị $x$ khay bánh mì và $y$ khay trái cây, với $x,y$ là các số nguyên không âm. Chi phí cho mỗi khay lần lượt là $40$ nghìn đồng và $30$ nghìn đồng; tổng chi phí không vượt quá $600$ nghìn đồng. Thời gian chuẩn bị mỗi khay lần lượt là $2$ phút và $3$ phút; tổng thời gian không vượt quá $42$ phút.
\begin{enumerate}
\item Tổng chi phí chuẩn bị là $40x+30y$ (nghìn đồng).
\item Phương án chuẩn bị $9$ khay bánh mì và $9$ khay trái cây đáp ứng cả hai giới hạn trên.
\item Các phương án đáp ứng yêu cầu được mô tả bởi
$\begin{cases}4x+3y\leq60,\\2x+3y\leq42,\\x,y\in\mathbb Z,\ x\geq0,\ y\geq0.\end{cases}$
\item Nếu chuẩn bị đúng $9$ khay bánh mì thì chuẩn bị được nhiều nhất $7$ khay trái cây.
\end{enumerate}
\item Cho hai tập hợp $A=[-2;4)$ và $B=(1;6]$.
\begin{enumerate}
\item $-2\in A$.
\item $\{4\}\subset A$.
\item $A\cap B=[1;4)$.
\item $B\setminus A=[4;6]$.
\end{enumerate}
\end{enumerate}
\vspace{7pt}
\parttitle{PHẦN III. Trả lời ngắn}
\textit{Ghi kết quả vào ô đáp số. Chỉ làm tròn khi câu hỏi yêu cầu.}
\begin{enumerate}
\item Cho $90^\circ<\alpha<180^\circ$ và $\cos\alpha=-\dfrac35$. Tính giá trị
\[
T=7\,\frac{\sin\alpha+2\cos\alpha}{\sin\alpha-\cos\alpha}.
\]
\answerbox
\item Cho $A=\{x\in\mathbb Z\mid -3\leq2x+1<12\}$ và $B=\{x\in\mathbb Z\mid x^2\leq9\}$. Tập hợp $A\cup B$ có bao nhiêu phần tử?
\par\answerbox
\item \begin{minipage}[t]{0.47\linewidth}
Hai vị trí $A,B$ trên mặt đất bằng phẳng cách nhau $40$ m, nằm cùng phía so với chân $H$ của một tòa nhà thẳng đứng; $A,B,H$ thẳng hàng theo thứ tự đó. Tại mỗi vị trí, người đo đặt máy có tâm ngắm cao $1{,}5$ m so với mặt đất. Góc nâng từ phương nằm ngang đến đỉnh tòa nhà đo tại $A$ là $30^\circ$, tại $B$ là $60^\circ$. Tòa nhà cao bao nhiêu mét? \textit{Làm tròn kết quả đến hàng phần mười.}
\par\vspace{4pt}\answerbox
\end{minipage}\hfill%
\begin{minipage}[t]{0.51\linewidth}
\vspace{-0.4\baselineskip}
\centering
\begin{tikzpicture}[scale=1.08, font=\normalsize, >=Stealth, line join=round, line cap=round]
    \draw[line width=1.2pt] (-0.2,0) -- (4.5,0) node[right=1pt, font=\small\bfseries] {mặt đất};
    \draw[line width=1.3pt, fill=gray!25] (3.8,0) rectangle (4.2,3.5);
    \coordinate (C) at (4.0,3.5);
    \node[above=2pt, font=\normalsize\bfseries] at (C) {$C$};
    \node[below=3pt, font=\normalsize\bfseries] at (4.0,0) {$H$};
    \draw[line width=1.2pt] (0.4,0) -- (0.4,0.6) coordinate (A1);
    \draw[line width=1.2pt] (2.4,0) -- (2.4,0.6) coordinate (B1);
    \fill (A1) circle (2.5pt) node[above left=-1pt, font=\normalsize\bfseries] {$A'$};
    \fill (B1) circle (2.5pt) node[above left=-1pt, font=\normalsize\bfseries] {$B'$};
    \fill (0.4,0) circle (2.5pt) node[below left=0pt, font=\normalsize\bfseries] {$A$};
    \fill (2.4,0) circle (2.5pt) node[below left=0pt, font=\normalsize\bfseries] {$B$};
    \draw[dashed, line width=0.9pt] (A1) -- (3.8,0.6) coordinate (H1);
    \draw[line width=1.5pt, ink] (A1) -- (C);
    \draw[line width=1.5pt, ink] (B1) -- (C);
    \draw[<->, line width=1.0pt] (0.4,-0.45) -- node[below=1pt, font=\small\bfseries] {$40\text{ m}$} (2.4,-0.45);
    \draw[<->, line width=1.0pt] (0.05,0) -- node[left=1pt, font=\small\bfseries] {$1{,}5\text{ m}$} (0.05,0.6);
    \pic[draw, line width=1.0pt, angle radius=8mm] {angle=H1--A1--C};
    \node at ($(A1)+(12:12mm)$) [font=\small\bfseries] {$30^\circ$};
    \pic[draw, line width=1.0pt, angle radius=7mm] {angle=H1--B1--C};
    \node at ($(B1)+(26:11mm)$) [font=\small\bfseries] {$60^\circ$};
\end{tikzpicture}
\end{minipage}
\item Gọi $D$ là tập các điểm $(x;y)$ thỏa mãn
\[
x\geq0,\quad y\geq0,\quad x+y\leq4,\quad 2x+y\leq6.
\]
Tìm giá trị lớn nhất của tham số thực $m$ để $(m-1)x+2y\leq8$ đúng với \textbf{mọi} điểm $(x;y)\in D$.
\par\answerbox
\end{enumerate}
\newpage
\parttitle{PHẦN IV. Tự luận}
\textit{Trình bày lập luận và các bước tính trên giấy làm bài.}
\begin{enumerate}
\item Cho hai tập hợp $A=(-\infty;3]$ và $B=[-1;5)$. Xác định $A\cap B$, $A\cup B$, $A\setminus B$; biểu diễn $A\setminus B$ trên trục số.
\item Cho $90^\circ<\alpha<180^\circ$ và
\[
P=\frac{1-\cos^2\alpha}{\sin\alpha\cos\alpha}.
\]
Chứng minh rằng $P=\tan\alpha$. Tính $P$ khi $\sin\alpha=\dfrac5{13}$.
\item Một thửa đất có dạng tứ giác lồi $ABCD$ như hình dưới đây. Người ta đo được $AB=80$ m, $BC=60$ m, $\widehat{ABC}=90^\circ$, $\widehat{DAC}=60^\circ$ và $\widehat{DCA}=45^\circ$. Tính diện tích thửa đất, \textit{làm tròn đến hàng đơn vị, đơn vị $\mathrm{m}^2$}.
\begin{center}
\begin{tikzpicture}[scale=1.35, line join=round, line cap=round, font=\normalsize]
\coordinate (A) at (0,0);
\coordinate (C) at (5,0);
\coordinate (B) at (3.2,-2.4);
\coordinate (D) at (1.830127,3.169873);
\fill[ink!8] (A)--(B)--(C)--(D)--cycle;
\draw[line width=1.4pt, ink] (A) -- (D) -- (C) -- (B) -- cycle;
\draw[dashed, line width=1.1pt, gray!80!black] (A) -- (C);
\draw[line width=1.4pt, ink] (A) -- node[below left=2pt, font=\normalsize\bfseries, black] {$80\ \mathrm{m}$} (B);
\draw[line width=1.4pt, ink] (B) -- node[below right=2pt, font=\normalsize\bfseries, black] {$60\ \mathrm{m}$} (C);
\pic[draw, line width=1.0pt, angle radius=8mm, angle eccentricity=1.45, "$60^\circ$", font=\normalsize\bfseries] {angle=C--A--D};
\pic[draw, line width=1.0pt, angle radius=9mm, angle eccentricity=1.45, "$45^\circ$", font=\normalsize\bfseries] {angle=D--C--A};
\pic[draw, line width=1.0pt, angle radius=5mm] {right angle=C--B--A};
\fill (A) circle (3pt) node[left=4pt, font=\Large\bfseries] {$A$};
\fill (B) circle (3pt) node[below=4pt, font=\Large\bfseries] {$B$};
\fill (C) circle (3pt) node[right=4pt, font=\Large\bfseries] {$C$};
\fill (D) circle (3pt) node[above=4pt, font=\Large\bfseries] {$D$};
\end{tikzpicture}
\end{center}
\item Một xưởng sản xuất hai loại sản phẩm $A$ và $B$ theo từng lô trọn vẹn. Mỗi lô cần lượng nguyên liệu và đem lại lợi nhuận như sau:
\begin{center}
\begin{tabular}{|l|c|c|}\hline
&\textbf{Một lô loại $A$}&\textbf{Một lô loại $B$}\\\hline
Nguyên liệu I&$2$ kg&$1$ kg\\\hline
Nguyên liệu II&$1$ kg&$2$ kg\\\hline
Lợi nhuận&$300\,000$ đồng&$250\,000$ đồng\\\hline
\end{tabular}
\end{center}
Xưởng có tối đa $20$ kg nguyên liệu I và $16$ kg nguyên liệu II. Giả sử sản phẩm làm ra đều bán hết, không có giới hạn nào khác.
\begin{enumerate}
\item Gọi $x,y$ lần lượt là số lô loại $A,B$. Lập hệ bất phương trình biểu thị các điều kiện của bài toán và viết biểu thức lợi nhuận.
\item Xưởng nên sản xuất bao nhiêu lô mỗi loại để lợi nhuận lớn nhất? Tính lợi nhuận lớn nhất đó.
\end{enumerate}
\end{enumerate}
\vfill
\begin{center}\textbf{--- HẾT ---}\end{center}
\newpage
\begin{center}
{\Large\bfseries BẢNG ĐÁP ÁN TỔNG HỢP}\\[5pt]
{\bfseries ÔN TẬP GIỮA KỲ I -- ĐỀ SỐ 01}
\end{center}
\parttitle{PHẦN I. Trắc nghiệm nhiều lựa chọn}
\begin{center}
\begin{tabularx}{\linewidth}{|>{\bfseries}l|*{12}{>{\centering\arraybackslash}X|}}\hline
Câu&1&2&3&4&5&6&7&8&9&10&11&12\\\hline
Đáp án&B&D&A&C&A&D&B&C&B&A&D&C\\\hline
\end{tabularx}
\end{center}
\parttitle{PHẦN II. Trắc nghiệm đúng/sai}
\begin{center}
\begin{tabular}{|c|c|c|c|c|}\hline
\textbf{Câu}&\textbf{a)}&\textbf{b)}&\textbf{c)}&\textbf{d)}\\\hline
1&Đúng&Sai&Đúng&Sai\\\hline
2&Đúng&Sai&Sai&Đúng\\\hline
\end{tabular}
\end{center}
\parttitle{PHẦN III. Trả lời ngắn}
\begin{center}
\begin{tabularx}{\linewidth}{|>{\bfseries}l|*{4}{>{\centering\arraybackslash}X|}}\hline
Câu&1&2&3&4\\\hline
Đáp số&$-2$&$9$&$36{,}1$ m&$3$\\\hline
\end{tabularx}
\end{center}
\parttitle{PHẦN IV. Tự luận -- kết quả cần đạt}
\begin{tabularx}{\linewidth}{|c|>{\raggedright\arraybackslash}X|}\hline
\textbf{Câu}&\textbf{Kết quả tóm tắt}\\\hline
1&$A\cap B=[-1;3]$; $A\cup B=(-\infty;5)$; $A\setminus B=(-\infty;-1)$. Trục số của $A\setminus B$: tia về bên trái, không lấy điểm $-1$.\\[4pt]\hline
2&$P=\tan\alpha$; khi $\sin\alpha=\dfrac5{13}$ thì $P=-\dfrac5{12}$.\\[4pt]\hline
3&Diện tích thửa đất $S=9900-2500\sqrt3\approx\boxed{5570\ \mathrm{m}^2}$.\\[4pt]\hline
4&Hệ BPT: $\begin{cases}2x+y\leq20\\x+2y\leq16\\x,y\in\mathbb Z,\ x\geq0,\ y\geq0\end{cases}$; Lợi nhuận: $L=300\,000x+250\,000y$ (đồng).\\
&Xưởng nên sản xuất \textbf{8 lô loại $A$} và \textbf{4 lô loại $B$}; lợi nhuận lớn nhất đạt $\boxed{3\,400\,000\text{ đồng}}$.\\\hline
\end{tabularx}

\newpage
\begin{center}
{\Large\bfseries HÌNH VẼ MIỀN NGHIỆM HỆ BẤT PHƯƠNG TRÌNH}\\[3pt]
{\normalsize\bfseries (PHẦN III -- CÂU 4 \& PHẦN IV -- CÂU 4)}
\end{center}

\vspace{6pt}
\begin{center}
\textbf{\large 1. Phần III -- Câu 4: Miền đa giác $D$ của hệ bất phương trình}\\[6pt]
\begin{tikzpicture}[scale=1.3, font=\normalsize, >=Stealth, line join=round, line cap=round]
    % Trục toạ độ
    \draw[->, line width=1.2pt] (-1,0) -- (5.6,0) node[below=2pt, font=\normalsize\bfseries] {$x$};
    \draw[->, line width=1.2pt] (0,-0.7) -- (0,7.2) node[left=2pt, font=\normalsize\bfseries] {$y$};
    \node[below left=2pt, font=\normalsize\bfseries] at (0,0) {$O$};
    
    % Miền D tô màu
    \fill[blue!15] (0,0) -- (3,0) -- (2,2) -- (0,4) -- cycle;
    
    % Đường thẳng d2: 2x+y=6 (kéo từ y=6.5 xuống đúng trục Ox tại x=3)
    \draw[line width=1.3pt, teal!80!black] (-0.25,6.5) -- (3,0);
    \node[above right, font=\normalsize\bfseries, teal!80!black] at (0.35,5.7) {$d_2: 2x+y=6$};
    
    % Đường thẳng d1: x+y=4 (kéo từ góc trên trái xuống đúng trục Ox tại x=4)
    \draw[line width=1.3pt, blue!80!black] (-0.8,4.8) -- (4,0);
    \node[above left, font=\normalsize\bfseries, blue!80!black] at (-0.8,4.8) {$d_1: x+y=4$};
    
    % Cạnh miền đa giác D (nét rất đậm)
    \draw[line width=2.0pt, ink] (0,0) -- (3,0) -- (2,2) -- (0,4) -- cycle;
    
    % Các đỉnh của miền D
    \fill (0,0) circle (3.5pt);
    \fill (3,0) circle (3.5pt);
    \node[below=4pt, font=\normalsize\bfseries, black] at (3,0) {$A(3;0)$};
    \fill (2,2) circle (3.8pt) node[above right=3pt, font=\normalsize\bfseries, black] {$B(2;2)$};
    \fill (0,4) circle (3.5pt) node[left=4pt, font=\normalsize\bfseries, black] {$C(0;4)$};
    
    % Các giao điểm trục khác
    \fill (0,6) circle (2.8pt) node[left=4pt, font=\normalsize\bfseries, black] {$6$};
    \fill (4,0) circle (2.8pt) node[below right=1pt, font=\normalsize\bfseries, black] {$4$};
    
    % Đường dóng đỉnh B(2;2)
    \draw[dashed, line width=1.0pt, gray!80!black] (2,0) -- (2,2) -- (0,2);
    \node[below left=1pt, font=\normalsize\bfseries, black] at (2,0) {$2$};
    \node[left=4pt, font=\normalsize\bfseries, black] at (0,2) {$2$};
    
    % Nhãn miền D
    \node[ink, font=\Large\bfseries] at (1.0, 1.2) {Miền $D$};
\end{tikzpicture}
\end{center}

\vspace{6pt}
\hrule
\vspace{6pt}

\begin{center}
\textbf{\large 2. Phần IV -- Câu 4: Miền nghiệm tứ giác $OABC$ của bài toán tối ưu}\\[6pt]
\begin{tikzpicture}[scale=0.8, font=\normalsize, >=Stealth, line join=round, line cap=round]
    % Trục toạ độ
    \draw[->, line width=1.2pt] (-1.5,0) -- (14.2,0) node[below=2pt, font=\normalsize\bfseries] {$x$ (lô $A$)};
    \draw[->, line width=1.2pt] (0,-1.0) -- (0,12.0) node[left=2pt, font=\normalsize\bfseries] {$y$ (lô $B$)};
    \node[below left=2pt, font=\normalsize\bfseries] at (0,0) {$O$};
    
    % Miền nghiệm tô màu
    \fill[ink!15] (0,0) -- (10,0) -- (8,4) -- (0,8) -- cycle;
    
    % Đường 2x+y=20 (kéo từ góc trên xuống đúng trục Ox tại x=10)
    \draw[line width=1.3pt, blue!80!black] (4.5,11) -- (10,0);
    \node[above right, font=\normalsize\bfseries, blue!80!black] at (4.6, 10.8) {$2x+y=20$};
    
    % Đường x+2y=16 (bắt đầu từ trục Oy tại y=8 kéo xuống dưới phải)
    \draw[line width=1.3pt, teal!80!black] (0,8) -- (14.5,0.75);
    \node[above right, font=\normalsize\bfseries, teal!80!black] at (14.5, 0.75) {$x+2y=16$};
    
    % Cạnh miền nghiệm tứ giác
    \draw[line width=2.0pt, ink] (0,0) -- (10,0) -- (8,4) -- (0,8) -- cycle;
    
    % Các đỉnh
    \fill (0,0) circle (3.8pt);
    \fill (10,0) circle (4.0pt);
    \node[below right=1pt, font=\normalsize\bfseries, black] at (10,0) {$A(10;0)$};
    \fill (8,4) circle (4.2pt) node[above right=3pt, font=\normalsize\bfseries, black] {$B(8;4)$};
    \fill (0,8) circle (4.0pt) node[left=4pt, font=\normalsize\bfseries, black] {$C(0;8)$};
    
    % Đường dóng đỉnh B(8;4)
    \draw[dashed, line width=1.0pt, gray!80!black] (8,0) -- (8,4) -- (0,4);
    \node[below left=1pt, font=\normalsize\bfseries, black] at (8,0) {$8$};
    \node[left=4pt, font=\normalsize\bfseries, black] at (0,4) {$4$};
    
    % Nhãn miền nghiệm
    \node[ink, font=\Large\bfseries] at (3.5, 2.6) {Miền tứ giác $OABC$};
\end{tikzpicture}
\end{center}
\end{document}
